% =====================================================================
\chapter{Gradiente de política}
\label{tema:2}
% =====================================================================

Este tema deduce el gradiente del rendimiento esperado respecto a los parámetros de la política y lo convierte en un algoritmo estimable por muestreo, REINFORCE~\cite{williams}. Después estudia la varianza del estimador y las dos técnicas que la reducen: la línea base y la ventaja estimada por un crítico, incluida la estimación generalizada de la ventaja (GAE)~\cite{gae}. Se siguen las §13.1--13.5 de Sutton y Barto~\cite{sb} y las clases~5 y~6 de CS285~\cite{cs285}.

\section{Políticas parametrizadas y función de puntuación}

\begin{definicion}[Rendimiento y función de puntuación]\label{def:score}
Sea $\pi_\theta$ una política parametrizada en un MDP episódico. El \emph{rendimiento} es $J(\theta)=\E_{\tau\sim p_\theta}[R(\tau)]$, con $R(\tau)=\sum_{t=0}^{T-1}\gamma^{t}r_{t+1}$ y $p_\theta=p_{\pi_\theta}$ dada por \eqref{eq:ptau}. Si $\pi_\theta(a\mid s)>0$, el vector $\nabla_\theta\log\pi_\theta(a\mid s)$ se llama \emph{función de puntuación} (\emph{score}) o vector de elegibilidad.
\end{definicion}

La parametrización más común en espacios de acciones discretos es la softmax de unas preferencias o logits~\cite[§13.1]{sb}. En un LLM los logits son la salida del Transformer y dependen de $\theta$ de forma no lineal; el caso tabular, en el que los propios logits son los parámetros, permite calcular todo a mano.

\begin{proposicion}[Puntuación de la softmax]\label{prop:score-softmax}
Sea $\pi_\theta=\softmax(\theta)$, con $\theta\in\R^{K}$ y $\pi_\theta(a)=e^{\theta_a}/\sum_b e^{\theta_b}$. Entonces
\[
\nabla_\theta\log\pi_\theta(a)=e_a-\pi_\theta .
\]
Si los logits son $z=f_\theta(s)$, por la regla de la cadena $\nabla_\theta\log\pi_\theta(a\mid s)=J_f(s)^{\top}(e_a-\pi_\theta(\cdot\mid s))$, con $J_f$ la matriz jacobiana de $f_\theta$.
\end{proposicion}

\begin{proof}
$\log\pi_\theta(a)=\theta_a-\log\sum_b e^{\theta_b}$. La derivada parcial respecto a $\theta_k$ es $\mathbb{1}[k=a]-e^{\theta_k}/\sum_b e^{\theta_b}=(e_a)_k-\pi_\theta(k)$. La segunda afirmación es la regla de la cadena del módulo~4.
\end{proof}

El vector $e_a-\pi_\theta$ es exactamente el gradiente, cambiado de signo, de la entropía cruzada con etiqueta $a$ (módulo~8). Esta coincidencia explica la relación entre gradiente de política y ajuste supervisado que se establece más abajo.

\section{El teorema del gradiente de política}

\begin{lema}[Truco de la derivada logarítmica]\label{lem:logder}
Si $p_\theta(z)>0$ es diferenciable en $\theta$, entonces $\nabla_\theta p_\theta(z)=p_\theta(z)\,\nabla_\theta\log p_\theta(z)$.
\end{lema}

\begin{proof}
Es la regla de la cadena aplicada a $\log$: $\nabla_\theta\log p_\theta(z)=\nabla_\theta p_\theta(z)/p_\theta(z)$.
\end{proof}

\begin{teorema}[Gradiente de política, forma de trayectorias]\label{teo:pg}
En un MDP episódico finito con política diferenciable y estrictamente positiva,
\[
\nabla_\theta J(\theta)=\E_{\tau\sim p_\theta}\Bigl[R(\tau)\sum_{t=0}^{T-1}\nabla_\theta\log\pi_\theta(a_t\mid s_t)\Bigr].
\]
\end{teorema}

\begin{proof}
Como el conjunto de trayectorias es finito, $J(\theta)=\sum_\tau p_\theta(\tau)R(\tau)$ y el gradiente puede introducirse en la suma. Por el \cref{lem:logder},
$\nabla_\theta J=\sum_\tau p_\theta(\tau)\,R(\tau)\,\nabla_\theta\log p_\theta(\tau)=\E_{p_\theta}[R(\tau)\nabla_\theta\log p_\theta(\tau)]$. Tomando logaritmos en \eqref{eq:ptau},
\[
\log p_\theta(\tau)=\log\rho_0(s_0)+\sum_{t}\log\pi_\theta(a_t\mid s_t)+\sum_t\log p(s_{t+1},r_{t+1}\mid s_t,a_t).
\]
El primer y el tercer sumando no dependen de $\theta$, luego su gradiente es nulo, y queda $\nabla_\theta\log p_\theta(\tau)=\sum_t\nabla_\theta\log\pi_\theta(a_t\mid s_t)$.
\end{proof}

Lo notable del teorema es que la dinámica $p$ desaparece del gradiente: basta con poder \emph{muestrear} trayectorias y evaluar su recompensa, aunque ni el entorno ni la recompensa sean derivables. En la forma de Sutton y Barto~\cite[§13.2]{sb}, el teorema se enuncia con la distribución de estados de la política y $Q^{\pi}$; la versión anterior es la de la clase~5 de CS285~\cite{cs285}.

\begin{corolario}[Gradiente de política en un LLM]\label{cor:pg-llm}
En el MDP de tokens (\cref{def:mdp-tokens}),
\[
\nabla_\theta J(\theta)=\E_{x}\,\E_{y\sim\pi_\theta(\cdot\mid x)}\Bigl[r(x,y)\sum_{t=1}^{T}\nabla_\theta\log\pi_\theta(y_t\mid x,y_{<t})\Bigr].
\]
\end{corolario}

\begin{proof}
Es el \cref{teo:pg} con $R(\tau)=r(x,y)$ (\cref{prop:mdp-tokens}), promediado sobre la distribución de instrucciones.
\end{proof}

El sumatorio $\sum_t\nabla_\theta\log\pi_\theta(y_t\mid x,y_{<t})$ es el gradiente de la log-verosimilitud de $y$, el mismo que maximiza el SFT del módulo~20. El gradiente de política es, por tanto, un SFT sobre las propias muestras del modelo, ponderado por su recompensa~\cite{cs336}.

\section{REINFORCE}

\begin{definicion}[Estimador REINFORCE]\label{def:reinforce}
Dadas $N$ trayectorias independientes $\tau^{(1)},\dots,\tau^{(N)}\sim p_\theta$, el estimador REINFORCE del gradiente es
\[
\widehat{g}=\frac1N\sum_{i=1}^{N}R(\tau^{(i)})\sum_{t}\nabla_\theta\log\pi_\theta(a^{(i)}_t\mid s^{(i)}_t),
\]
y el algoritmo actualiza $\theta\leftarrow\theta+\eta\,\widehat{g}$ con una tasa de aprendizaje $\eta>0$.
\end{definicion}

Por el \cref{teo:pg}, $\E[\widehat g]=\nabla_\theta J$: el estimador es insesgado. Williams~\cite{williams} lo introdujo en 1992 como un caso de una familia de algoritmos que siguen el gradiente en esperanza.

\begin{ejemplo}[Un paso de REINFORCE en un bandido]\label{ej:reinforce}
Sean tres acciones con logits $\theta=(0,0,0)$ y recompensas $r=(1,2,0)$. Entonces $\pi=(1/3,1/3,1/3)$ y $J=\pi\cdot r=1$. El gradiente exacto es $\partial J/\partial\theta_k=\sum_a\pi(a)r(a)(e_a-\pi)_k=\pi_k(r_k-J)$, es decir, $\nabla J=(0;\ 0{,}333;\ -0{,}333)$. Si se muestrea $a=1$ (con $r=2$), la puntuación es $e_1-\pi=(-0{,}333;\ 0{,}667;\ -0{,}333)$ y la estimación es $2\,(e_1-\pi)=(-0{,}667;\ 1{,}333;\ -0{,}667)$. Con $\eta=1$, $\theta=(-0{,}667;\ 1{,}333;\ -0{,}667)$ y $\pi=\softmax(\theta)=(0{,}107;\ 0{,}787;\ 0{,}107)$.
\end{ejemplo}

\section{Causalidad y recompensa futura}

\begin{proposicion}[Las recompensas pasadas no aportan]\label{prop:causalidad}
Si $t'<t$, entonces $\E_{p_\theta}\bigl[r_{t'+1}\,\nabla_\theta\log\pi_\theta(a_t\mid s_t)\bigr]=0$. En consecuencia,
\[
\nabla_\theta J(\theta)=\E_{p_\theta}\Bigl[\sum_{t}\nabla_\theta\log\pi_\theta(a_t\mid s_t)\,\widehat{Q}_t\Bigr],\qquad \widehat Q_t=\sum_{t'\ge t}\gamma^{t'}r_{t'+1}.
\]
\end{proposicion}

\begin{proof}
Condicionando a la historia $h_t=(s_0,a_0,\dots,s_t)$, que determina $r_{t'+1}$ para $t'<t$, se tiene
$\E[r_{t'+1}\nabla\log\pi_\theta(a_t\mid s_t)\mid h_t]=r_{t'+1}\sum_{a}\pi_\theta(a\mid s_t)\nabla_\theta\log\pi_\theta(a\mid s_t)=r_{t'+1}\nabla_\theta\sum_a\pi_\theta(a\mid s_t)=r_{t'+1}\nabla_\theta 1=0$, por el \cref{lem:logder}. La ley de la esperanza total da el primer enunciado. Para el segundo se desarrolla $R(\tau)=\sum_{t'}\gamma^{t'}r_{t'+1}$ en el \cref{teo:pg} y se eliminan los términos con $t'<t$.
\end{proof}

La cantidad $\widehat Q_t$ (la «recompensa futura», \emph{reward-to-go}) es una estimación de una muestra de $Q^{\pi}(s_t,a_t)$, salvo el factor $\gamma^{t}$~\cite{cs285hw2}. En el MDP de tokens, con una sola recompensa final, todas las $\widehat Q_t$ valen $r(x,y)$ y la proposición no aporta nada; su utilidad aparece con recompensas intermedias, como la penalización KL por token del \cref{tema:3}.

\section{Línea base y varianza}

\begin{lema}[La puntuación tiene media nula]\label{lem:score-media}
Para todo estado $s$, $\E_{a\sim\pi_\theta(\cdot\mid s)}\bigl[\nabla_\theta\log\pi_\theta(a\mid s)\bigr]=0$.
\end{lema}

\begin{proof}
$\sum_a\pi_\theta(a\mid s)\,\nabla_\theta\pi_\theta(a\mid s)/\pi_\theta(a\mid s)=\nabla_\theta\sum_a\pi_\theta(a\mid s)=\nabla_\theta 1=0$.
\end{proof}

\begin{proposicion}[Línea base]\label{prop:baseline}
Sea $b(s)$ cualquier función del estado que no dependa de la acción. Entonces
\[
\nabla_\theta J(\theta)=\E_{p_\theta}\Bigl[\sum_t\nabla_\theta\log\pi_\theta(a_t\mid s_t)\bigl(\widehat Q_t-b(s_t)\bigr)\Bigr].
\]
\end{proposicion}

\begin{proof}
Basta ver que $\E[\nabla_\theta\log\pi_\theta(a_t\mid s_t)\,b(s_t)]=0$. Condicionando a $s_t$, el factor $b(s_t)$ sale de la esperanza y queda $b(s_t)\cdot 0$ por el \cref{lem:score-media}.
\end{proof}

La línea base no cambia la esperanza, pero sí la varianza. En el \cref{ej:reinforce}, la varianza total (traza de la matriz de covarianzas) del estimador de una muestra $(r(a)-b)(e_a-\pi)$ vale $0{,}889$ con $b=0$ y $0{,}222$ con $b=J=1$. Si se suma~10 a todas las recompensas, con $b=0$ sube a $80{,}9$, mientras que con $b=J$ sigue siendo $0{,}222$: la línea base hace el método insensible al origen arbitrario de la escala de recompensas. La clase~5 de CS285 observa que la recompensa media no es la línea base óptima, aunque sí una muy buena~\cite{cs285}; la proposición siguiente da la óptima.

\begin{proposicion}[Línea base de varianza mínima]\label{prop:baseline-optima}
En un bandido con $g(a)=\nabla_\theta\log\pi_\theta(a)$, la varianza total del estimador $(r(a)-b)\,g(a)$ con $a\sim\pi_\theta$ es mínima para
\[
b^{*}=\frac{\E\bigl[r(a)\,\norm{g(a)}^{2}\bigr]}{\E\bigl[\norm{g(a)}^{2}\bigr]}.
\]
\end{proposicion}

\begin{proof}
La varianza total es $V(b)=\E[(r-b)^{2}\norm{g}^{2}]-\norm{\E[(r-b)g]}^{2}$. Por el \cref{lem:score-media}, $\E[(r-b)g]=\E[rg]=\nabla J$ no depende de $b$. Así, $V(b)$ es un polinomio de grado~2 en $b$ con $V'(b)=-2\E[r\norm{g}^{2}]+2b\,\E[\norm{g}^{2}]$ y $V''(b)=2\E[\norm g^{2}]>0$ (salvo en el caso trivial $g\equiv 0$). Su único punto crítico, $b^{*}$, es el mínimo.
\end{proof}

\begin{ejemplo}\label{ej:baseline-optima}
Con $\pi=(0{,}5;\ 0{,}3;\ 0{,}2)$ y $r=(1,2,0)$, las normas $\norm{e_a-\pi}^{2}$ valen $0{,}38$, $0{,}78$ y $0{,}98$. Entonces $b^{*}=0{,}658/0{,}62=1{,}061$, frente a $\E[r]=1{,}1$. Las varianzas son $1{,}002$ con $b=0$, $0{,}305$ con $b=\E[r]$ y $0{,}304$ con $b^{*}$: la recompensa media es casi óptima y mucho más sencilla de estimar.
\end{ejemplo}

\section{Ventaja, crítico y estimación generalizada de la ventaja}

La línea base natural en un MDP es $b(s_t)=V^{\pi}(s_t)$. Entonces $\widehat Q_t-V^{\pi}(s_t)$ es una estimación de la ventaja $A^{\pi}(s_t,a_t)$ de la \cref{def:valor}, y el gradiente queda $\E[\sum_t\nabla\log\pi_\theta(a_t\mid s_t)\,A^{\pi}(s_t,a_t)]$, salvo los factores de descuento. Como $V^{\pi}$ es desconocida, se aproxima con una red $V_\phi$, el \emph{crítico}, entrenada por regresión sobre los retornos observados. La política es el \emph{actor}: de ahí el nombre de métodos actor-crítico~\cite[§13.5]{sb}.

\begin{definicion}[Error de diferencia temporal y GAE]\label{def:gae}
Dado un crítico $V$, el \emph{error de diferencia temporal} es $\delta_t=r_{t+1}+\gamma V(s_{t+1})-V(s_t)$, con $V(s_T)=0$ en el estado terminal. Para $\lambda\in[0,1]$, la \emph{estimación generalizada de la ventaja} es~\cite{gae}
\[
\widehat A^{\mathrm{GAE}(\gamma,\lambda)}_t=\sum_{l=0}^{T-t-1}(\gamma\lambda)^{l}\,\delta_{t+l},
\]
que satisface la recursión $\widehat A_t=\delta_t+\gamma\lambda\,\widehat A_{t+1}$, con $\widehat A_T=0$.
\end{definicion}

\begin{proposicion}[Casos extremos de GAE]\label{prop:gae}
Con $\lambda=0$, $\widehat A_t=\delta_t$. Con $\lambda=1$, $\widehat A_t=\sum_{l\ge0}\gamma^{l}r_{t+l+1}-V(s_t)$, el retorno observado menos la línea base. Si $V=V^{\pi}$, entonces $\E[\delta_t\mid s_t,a_t]=A^{\pi}(s_t,a_t)$.
\end{proposicion}

\begin{proof}
El caso $\lambda=0$ es inmediato. Con $\lambda=1$, la suma es telescópica:
$\sum_{l}\gamma^{l}\bigl[r_{t+l+1}+\gamma V(s_{t+l+1})-V(s_{t+l})\bigr]=\sum_l\gamma^{l}r_{t+l+1}+\sum_{l}\bigl[\gamma^{l+1}V(s_{t+l+1})-\gamma^{l}V(s_{t+l})\bigr]$, y la segunda suma se reduce a $\gamma^{T-t}V(s_T)-V(s_t)=-V(s_t)$. Por último, $\E[r_{t+1}+\gamma V^{\pi}(s_{t+1})\mid s_t,a_t]=Q^{\pi}(s_t,a_t)$ por el \cref{teo:bellman}, luego $\E[\delta_t\mid s_t,a_t]=Q^{\pi}-V^{\pi}=A^{\pi}$.
\end{proof}

El parámetro $\lambda$ regula el compromiso entre sesgo y varianza: con $\lambda=0$ se usa solo un paso de recompensa real y se confía en el crítico (poca varianza, sesgo si $V$ es inexacto); con $\lambda=1$ se usa el retorno completo (sin sesgo, mucha varianza)~\cite{cs285}. PPO calcula sus ventajas con GAE.

\begin{ejemplo}[Cálculo de GAE]\label{ej:gae}
Sea un episodio de tres pasos con recompensas $(0,0,1)$, valores del crítico $V(s_0)=0{,}5$, $V(s_1)=0{,}6$, $V(s_2)=0{,}8$, $\gamma=1$ y $\lambda=0{,}5$. Los errores son $\delta_0=0+0{,}6-0{,}5=0{,}1$, $\delta_1=0+0{,}8-0{,}6=0{,}2$ y $\delta_2=1+0-0{,}8=0{,}2$. Hacia atrás: $\widehat A_2=0{,}2$, $\widehat A_1=0{,}2+0{,}5\cdot 0{,}2=0{,}3$ y $\widehat A_0=0{,}1+0{,}5\cdot 0{,}3=0{,}25$. Con $\lambda=1$ se obtendría $\widehat A_0=0{,}1+0{,}2+0{,}2=0{,}5=G_0-V(s_0)$, de acuerdo con la \cref{prop:gae}.
\end{ejemplo}

\begin{lecturas}
\item Sutton y Barto~\cite{sb}, cap.~13: §13.1 (aproximación de la política), §13.2 «The Policy Gradient Theorem», §13.3 «REINFORCE: Monte Carlo Policy Gradient», §13.4 «REINFORCE with Baseline» y §13.5 «Actor--Critic Methods».
\item Levine, CS285~\cite{cs285}, clases~5 «Policy Gradients» y~6 «Actor-Critic Algorithms».
\item R.~J. Williams~\cite{williams}, artículo original de REINFORCE.
\item J.~Schulman \emph{et al.}~\cite{gae}, artículo original de GAE.
\end{lecturas}

\section*{Ejercicios del tema 2}
\addcontentsline{toc}{section}{Ejercicios del tema 2}

\begin{ejercicio}{1}{eje:2-baseline}
Explique por qué restar una línea base $b$ no sesga el gradiente de política pero sí cambia su varianza. ¿Por qué no puede usarse una línea base que dependa de la acción muestreada?
\end{ejercicio}

\begin{ejercicio}{2}{eje:2-bernoulli}
Sea la política de Bernoulli $\pi_\theta(1)=\sigma(\theta)$, $\pi_\theta(0)=1-\sigma(\theta)$, con $\theta=\ln 3$, y recompensas $r(1)=1$, $r(0)=0$. Calcule $\pi_\theta(1)$, $\nabla_\theta\log\pi_\theta(1)$, $\nabla_\theta\log\pi_\theta(0)$ y $\nabla_\theta J$, y compruebe que coincide con $\E[r\,\nabla\log\pi_\theta]$.
\end{ejercicio}

\begin{ejercicio}{2}{eje:2-paso}
Tres acciones con $\theta=(0,0,0)$. Se muestrea $a=2$ y se recibe $r=-1$. Aplique un paso de REINFORCE con $\eta=1$ y sin línea base. ¿Qué $\theta$ y qué $\pi$ se obtienen?
\end{ejercicio}

\begin{ejercicio}{2}{eje:2-cs285}
En el MDP del \cref{eje:1-cs285}, con $\gamma=1$, una trayectoria que elige $a_1$ exactamente $k$ veces antes de $a_2$ tiene probabilidad $\theta^{k}(1-\theta)$ y retorno $R=k$. (a) Calcule $\nabla_\theta\log p_\theta(\tau)$ en función de $k$. (b) Calcule $\nabla_\theta J$ con el \cref{teo:pg}, usando $\E[k]=\theta/(1-\theta)$ y $\E[k^{2}]=\theta(1+\theta)/(1-\theta)^{2}$. (c) Compruebe el resultado derivando $J(\theta)=\theta/(1-\theta)$ y evalúe en $\theta=1/2$.
\fuente{adaptado de Levine, CS285, otoño de 2023, tarea~2 «Policy Gradients», §8, problema~1(a)~\cite{cs285hw2}.}
\end{ejercicio}

\begin{ejercicio}{2}{eje:2-gae}
Un episodio de tres pasos tiene recompensas $(1,0,2)$ y valores del crítico $V(s_0)=2$, $V(s_1)=1{,}5$, $V(s_2)=1$, con $V(s_3)=0$. Con $\gamma=0{,}9$ y $\lambda=0{,}8$, calcule los errores $\delta_t$ y las ventajas GAE. Compruebe que con $\lambda=1$ se obtiene $\widehat A_0=G_0-V(s_0)$.
\end{ejercicio}

\begin{ejercicio}{3}{eje:2-optima}
Demuestre que $\E_{a\sim\pi}[\nabla_\theta\log\pi_\theta(a)]=0$, deduzca que una línea base constante no introduce sesgo y demuestre la fórmula de $b^{*}$ de la \cref{prop:baseline-optima}. Compruebe numéricamente el valor del \cref{ej:baseline-optima}.
\end{ejercicio}

\begin{ejercicio}{3}{eje:2-varianza}
En el MDP del \cref{eje:2-cs285}, el estimador de una trayectoria es $g(k)=k\bigl(k/\theta-1/(1-\theta)\bigr)$. Sabiendo que $\E[k^{3}]=\theta(\theta^{2}+4\theta+1)/(1-\theta)^{3}$ y $\E[k^{4}]=\theta(1+\theta)(\theta^{2}+10\theta+1)/(1-\theta)^{4}$, demuestre que
\[
\Var[g]=\frac{4\theta^{2}+8\theta+1}{\theta(1-\theta)^{4}},
\]
evalúela en $\theta=1/2$ y estudie su comportamiento cuando $\theta\to 0$ y $\theta\to 1$.
\fuente{adaptado de Levine, CS285, otoño de 2023, tarea~2 «Policy Gradients», §8, problema~2~\cite{cs285hw2}.}
\end{ejercicio}
